
\documentclass[a4paper,12pt]{article}
\usepackage{amsmath}
\usepackage{graphicx}
\usepackage{multicol}
\usepackage{geometry}
\usepackage{fancyhdr}
\usepackage{tikz}
\usepackage{eso-pic}

\geometry{top=1.5cm, bottom=1.5cm, left=1.5cm, right=1.5cm}

\pagestyle{fancy}
\fancyhf{}
\rhead{Esc. 4-124 Reynaldo Merín}
\lhead{Electrotecnia II}
\cfoot{\thepage}

\AddToShipoutPicture*{
    \put(150,550){\rotatebox{45}{\textcolor[gray]{0.95}{EXAMEN ORIGINAL NO DIGITALIZAR}}}
}

\begin{document}

\begin{center}
    \Large \textbf{Evaluación B – Electrotecnia II} \\
    \normalsize Duración: 60 minutos \\
    Alumno/a: \underline{\hspace{6cm}} \hspace{0.5cm} Fecha: \underline{\hspace{2cm}}
\end{center}

\vspace{0.3cm}

\begin{multicols}{2}

\section*{Problemas}

\begin{enumerate}
\item Un conductor de $l = 0{,}3\,m$ con $I = 8\,A$ está inmerso en un campo de $B = 1{,}1\,T$. Calcular la fuerza magnética.

\item Hallar la inducción $B$ si un conductor de $l = 15\,cm$ e intensidad $I = 18\,A$ recibe una fuerza de $3{,}5\,N$.

\item Un solenoide de $N = 1000$ espiras, $r = 0{,}1\,m$, $l = 0{,}2\,m$ y núcleo de hierro con $\mu = 2{,}5 \times 10^{-6}$. Calcular $L$.

\item Una bobina tiene $L = 0{,}005\,H$ y $\frac{di}{dt} = 0{,}6\,A/s$. Calcular la F.E.M. inducida.

\item Una bobina de $N = 1200$, $L = 0{,}8\,H$, $r = 0{,}04\,m$, $l = 0{,}15\,m$. Calcular $\mu$ del núcleo.

\item Un conductor de $l = 22\,cm$ con $I = 5\,A$ está perpendicular a un campo de $B = 0{,}6\,T$. Calcular la fuerza.

\item Calcular la corriente necesaria para que una bobina de $L = 0{,}02\,H$ tenga una F.E.M. inducida de $0{,}3\,V$ si $\frac{di}{dt}$ es constante.

\item Justificar con regla de la mano izquierda la dirección de la fuerza en un conductor con corriente hacia el sur en un campo de oeste a este.

\end{enumerate}

\end{multicols}

\end{document}